\documentclass[12pt,a4paper]{ctexart}

% ==================== 基础排版 ====================
\usepackage[margin=2.25cm,headheight=14pt]{geometry}
\usepackage{enumitem}
\usepackage[most]{tcolorbox}
\usepackage{xcolor}
\usepackage{tikz}
\usepackage{fancyhdr}
\usepackage{bookmark}
\usepackage{hyperref}
\hypersetup{hidelinks}

% ==================== 列表与段落 ====================
\setlist[itemize]{leftmargin=2.2em,itemsep=0.4em,topsep=0.4em}
\setlist[enumerate]{leftmargin=2.4em,itemsep=0.4em,topsep=0.4em}
\setlength{\parindent}{2em}
\linespread{1.18}

% ==================== 色彩定义 ====================
\definecolor{KeyYellow}     {RGB}{255,247,177}
\definecolor{KeyYellowFrame}{RGB}{210,175,0}
\definecolor{KeyGreen}      {RGB}{220,239,205}
\definecolor{KeyGreenFrame} {RGB}{69,148,77}
\definecolor{KeyLineGreen}  {RGB}{88,190,88}
\definecolor{AccentBlue}    {RGB}{30,100,180}
\definecolor{LightGray}     {RGB}{245,245,245}
\definecolor{MidGray}       {RGB}{180,180,180}

% ==================== 页眉页脚 ====================
\fancypagestyle{plain}{}
\pagestyle{fancy}
\fancyhf{}
\fancyhead[L]{\small\color{MidGray}\textsf{习概·期末复习}}
\fancyhead[R]{\small\color{MidGray}\textsf{\leftmark}}
\fancyfoot[C]{\small\color{MidGray}—\ \thepage\ —}
\renewcommand{\headrulewidth}{0.4pt}
\renewcommand{\footrulewidth}{0pt}
\fancypagestyle{plain}{
  \fancyhf{}
  \fancyfoot[C]{\small\color{MidGray}—\ \thepage\ —}
  \renewcommand{\headrulewidth}{0pt}
}

% ==================== 节标题样式 ====================
\ctexset{
  section = {
    format     = \Large\bfseries\color{AccentBlue},
    aftername  = \quad,
    afterskip  = 1.2ex plus 0.3ex minus 0.2ex,
  },
  subsection = {
    format     = \large\bfseries\color{AccentBlue!80!black},
    aftername  = \quad,
    afterskip  = 1.0ex plus 0.2ex minus 0.15ex,
  },
}

% ==================== 彩色信息盒 ====================
\newtcolorbox{yellowbox}[1]{
  enhanced,
  breakable,
  colback      = white,
  colbacktitle = white,
  colframe     = KeyYellowFrame,
  coltitle     = black,
  fonttitle    = \bfseries\small,
  title        = {#1},
  boxrule      = 0.8pt,
  arc          = 2mm,
  outer arc    = 2mm,
  left         = 3mm,
  right        = 3mm,
  top          = 2mm,
  bottom       = 2mm,
  toptitle     = 1.8mm,
  bottomtitle  = 1.8mm,
  drop fuzzy shadow = {color=black!15},
  before skip  = 1.2ex,
  after skip   = 1.5ex,
}

\newtcolorbox{greenbox}[1]{
  enhanced,
  breakable,
  colback      = white,
  colbacktitle = white,
  colframe     = KeyGreenFrame,
  coltitle     = black,
  fonttitle    = \bfseries\small,
  title        = {#1},
  boxrule      = 0.8pt,
  arc          = 2mm,
  outer arc    = 2mm,
  left         = 3mm,
  right        = 3mm,
  top          = 2mm,
  bottom       = 2mm,
  toptitle     = 1.8mm,
  bottomtitle  = 1.8mm,
  drop fuzzy shadow = {color=black!15},
  before skip  = 1.2ex,
  after skip   = 1.5ex,
}

\newtcolorbox{greenlinebox}[1]{
  enhanced,
  breakable,
  colback      = white,
  colframe     = KeyLineGreen,
  coltitle     = black,
  colbacktitle = white,
  fonttitle    = \bfseries\small,
  title        = {#1},
  boxrule      = 1.1pt,
  arc          = 2mm,
  outer arc    = 2mm,
  left         = 3mm,
  right        = 3mm,
  top          = 2mm,
  bottom       = 2mm,
  toptitle     = 1.8mm,
  bottomtitle  = 1.8mm,
  drop fuzzy shadow = {color=black!12},
  before skip  = 1.2ex,
  after skip   = 1.5ex,
}

% ==================== 标题页 ====================
\title{}
\author{}
\date{}

\begin{document}



% ---------- 标记说明 ----------


% ================================================================
\section{黄色重点：三个专业高度重合}
% ================================================================

\subsection{全面深化改革与高质量发展}

\begin{yellowbox}{全面深化改革要坚持正确方法论}
\begin{enumerate}
  \item 增强全面深化改革的系统性、整体性、协同性；
  \item 加强顶层设计和摸着石头过河相结合；
  \item 统筹改革发展稳定；
  \item 胆子要大，步子要稳；
  \item 坚持重大改革于法有据。
\end{enumerate}
\end{yellowbox}



\begin{yellowbox}{基本经济制度与民营经济}
\begin{enumerate}
  \item 民营经济是非公有制经济的重要经济组织形式，是推进中国式现代化的生力军，是高质量发展的重要基础，是推动我国全面建成社会主义现代化强国、实现第二个百年奋斗目标的重要力量。
\end{enumerate}
\end{yellowbox}


\subsection{全过程人民民主与全面依法治国}

\begin{yellowbox}{全过程人民民主}
\begin{enumerate}
  \item 全过程人民民主是全链条、全方位、全覆盖的民主。
  \item 全过程人民民主是最广泛、最真实、最管用的民主。
  \item 协商民主是实践全过程人民民主的重要形式。
  \item 基层民主是全过程人民民主的重要体现。
  \item 统战工作的本质要求是大团结、大联合，解决的就是人心和力量问题。
\end{enumerate}
\end{yellowbox}

\begin{yellowbox}{习近平法治思想与全面依法治国}
\begin{enumerate}
  \item 习近平法治思想的主要内容集中体现为“十二个坚持”；重点记忆第十二个：依法治国和依规治党的有机统一。
  \item 统筹处理全面依法治国的重大关系：
  \begin{enumerate}
    \item 正确处理政治和法治的关系；
    \item 正确处理改革和法治的关系；
    \item 正确处理依法治国和以德治国的关系；
    \item 正确处理依法治国和依规治党的关系。
  \end{enumerate}
\end{enumerate}
\end{yellowbox}



\begin{yellowbox}{“一国两制”的科学内涵}
\begin{enumerate}
  \item 牢牢把握“一国两制”的根本宗旨：坚定不移，确保不会变、不动摇；全面准确，确保不走样、不变形。
  \item 准确把握“一国”和“两制”的关系：一国是实行两制的前提和基础；两制从属和派生于一国，并统一于一国之内。
  \item 坚持中央全面管治权和保障特别行政区高度自治权统一。
  \item 坚定落实爱国者治港、爱国者治澳原则。
\end{enumerate}
\end{yellowbox}


% ================================================================
\section{绿色画线：三个专业重合，可能作为大题}
% ================================================================

\subsection{牢固树立和践行正确的政绩观}

\begin{greenlinebox}{大题要点：正确政绩观}
中国共产党把为民办事、为民造福作为最重要的政绩，把为老百姓办了多少好事实事作为检验政绩的重要标准。

评价干部的工作成效，不仅要看纸面上的指标数据，更要看人民的生活好不好、民生是否改善、生态是否良好、社会是否进步；不能简单以国内生产总值增长率论英雄，更不能搞脱离实际的盲目攀比，不搞劳民伤财的形象工程和政绩工程。

要以“功成不必在我”的精神境界和“功成必定有我”的历史担当，多做为后人作铺垫、打基础、利长远的好事，追求人民群众的好口碑、历史沉淀之后的真评价。
\end{greenlinebox}

\subsection{全面推进依法治国的总抓手}

\begin{greenlinebox}{大题要点：建设中国特色社会主义法治体系}
必须加快形成完备的法律规范体系、高效的法治实施体系、严密的法治监督体系、有力的法治保障体系，形成完善的党内法规体系。
\begin{enumerate}
  \item \textbf{完备的法律规范体系：}以宪法为核心，由部门齐全、结构严谨、内部协调、体例科学、调整有效的法律规范构成有机整体。
  \item \textbf{高效的法治实施体系：}包括执法、司法、守法等各个环节协调高效运转；把平等保护贯彻到立法、执法、司法、守法等各个环节，依法平等保护各类市场主体产权和合法权益。
  \item \textbf{严密的法治监督体系：}以规范和约束公权力运行为重点；完善权力运行制约和监督机制，规范立法、执法、司法机关权力行使，构建党的统一领导、全面覆盖、权威高效的法治监督体系。
  \item \textbf{有力的法治保障体系：}包括政治、思想、组织、制度、队伍、科技等保障条件；充分运用大数据、云计算、人工智能等现代科技手段，全面建设智慧法治，推进法治中国建设的数据化、网络化、智能化。
  \item \textbf{完善的党内法规体系：}把党内法规体系纳入中国特色社会主义法治体系，是我国法治区别于其他国家法治的鲜明特征；全面依法治国必须构建内容科学、程序严密、配套完备、运行有效的党内法规体系。
\end{enumerate}
\end{greenlinebox}

\subsection{文化繁荣兴盛与社会主义意识形态}

\begin{greenlinebox}{大题要点：文化繁荣兴盛的重要性}
\begin{enumerate}
  \item 文化繁荣兴盛是实现中华民族伟大复兴的精神支撑；
  \item 文化繁荣兴盛是建设社会主义现代化强国的应有之义；
  \item 文化繁荣兴盛是满足人民日益增长的美好生活需要的内在要求；
  \item 文化繁荣兴盛是在世界文化激荡中站稳脚跟的前提基础。
\end{enumerate}
\end{greenlinebox}

\begin{greenlinebox}{大题要点：坚持马克思主义在意识形态领域指导地位的根本制度}
\begin{enumerate}
  \item 坚持马克思主义在意识形态领域指导地位的制度，是中国特色社会主义制度体系的一项根本制度。
  \item 坚持和加强党对宣传思想文化工作全面领导，是发展社会主义先进文化的有力保障。
  \item 坚持马克思主义在意识形态领域指导地位的根本制度是历史的结论。马克思主义深刻揭示人类社会发展规律，是来自人民、为了人民、造福人民的思想体系，是指引人民认识世界、改造世界的行动指南。
  \item 坚持马克思主义在意识形态领域指导地位的根本制度，是坚持和巩固我国社会主义制度、保障我国文化建设正确方向的必然要求。
  \item 面对世界范围的思想文化相互激荡，我国社会思想观念和价值取向日趋多样。
  \item 只有坚持以马克思主义统领多样化思想文化发展，才能牢牢把握社会主义先进文化前进方向，夯实共同的思想基础，拉紧共同的精神纽带，促进全体人民在思想上、精神上凝聚在一起。
\end{enumerate}
\end{greenlinebox}

\subsection{推动构建新型国际关系}

\begin{greenlinebox}{大题要点：新型国际关系的主要内容}
\begin{enumerate}
  \item 坚持在和平共处五项原则基础上同各国发展友好合作，深化拓展平等、开放、合作的全球伙伴关系。
  \item 这是新时代中国外交理论和实践的重要创新。
  \item 推进大国协调合作，打造周边命运共同体，加强同发展中国家团结合作，不断完善我国全方位、多层次、立体化的外交布局，构建新型政党关系，开创国际关系新模式。
  \item 推进大国协调合作，构建和平共处、总体稳定、均衡发展的大国关系格局，是中国对外关系的重要内容。
  \item 秉持真实亲诚理念和正确义利观，加强同发展中国家团结合作。
\end{enumerate}
\end{greenlinebox}
\end{document}
